\documentclass[12pt]{article}
\usepackage[utf8]{inputenc}
\usepackage[T1]{fontenc}
\usepackage[margin=1in]{geometry}
\usepackage{setspace}
\usepackage{xcolor}
\usepackage[colorlinks=true, urlcolor=blue!60!black, linkcolor=blue!60!black]{hyperref}
\setlength{\parindent}{0em}
\setlength{\parskip}{0.6em}
\renewcommand{\baselinestretch}{1.0}
\begin{document}
\begin{center}
{\LARGE \textit{The Slaughter-House Cases}, 83 U.S. 36 (1873)\par}
\vspace{0.5em}
{\large Case Brief --- Two-Page Edition\par}
\end{center}

\section{Memory Jogger}
In its first reading of the Reconstruction Amendments, the Court upheld a Louisiana slaughterhouse monopoly and confined the Privileges or Immunities Clause to rights owing their existence to the federal government (\href{https://www.courtlistener.com/opinion/88661/butchers-benevolent-assn-v-crescent-city-live-stock-landing/}{83 U.S. 36}).

\section{Detailed Case Facts}
To stop slaughterhouse waste from fouling New Orleans's water supply, Louisiana in 1869 gave the Crescent City Company a 25-year exclusive right to land and slaughter livestock in three parishes, closing competitors but requiring the company to admit any butcher at capped fees. About 1,000 displaced butchers, represented by former Justice John A. Campbell, challenged the monopoly under the Thirteenth and Fourteenth Amendments.

\section{Procedural History}
The Louisiana Supreme Court upheld the statute; three consolidated cases reached the Supreme Court on writs of error, which affirmed after reargument in 1873.

\section{Judicial Votes}
Affirmed, 5--4. Miller wrote for Clifford, Davis, Strong, and Hunt; Field dissented, joined by Chase, Swayne, and Bradley, and Bradley and Swayne also wrote separately.

\section{Holding}
The Privileges or Immunities Clause protects only rights of national citizenship, not ordinary civil and economic rights of state citizenship such as pursuing a trade; the monopoly imposed no involuntary servitude and violated neither due process nor equal protection.

\section{Analysis of Opinions}
Miller treated the law as a police-power health measure and read the Amendments through their ``one pervading purpose,'' the freedom of the formerly enslaved. Because Section 1 distinguishes national from state citizenship and protects only national privileges, a broader reading would make the Court a ``perpetual censor upon all legislation of the States.'' Field called the health rationale a ``shallow'' pretence and argued the majority made the Clause ``a vain and idle enactment, which accomplished nothing''; it protects the right to pursue a lawful employment. Bradley would protect the right to choose a calling and the Bill of Rights, and Swayne answered that ``the novelty was known and the measure deliberately adopted.''

\section{Examples: Future Applications}
\textbf{Same-Side:} a single regulated waste-hauling franchise survives a displaced hauler's Privileges or Immunities claim. \textbf{Opposite-Side:} a head tax on travel to Washington, D.C., abridges a privilege Miller recognized. \textbf{Fence-Sitter:} incorporating the Bill of Rights through the Clause, as Justices Thomas and Gorsuch have urged.

\section{Critique}
Originalists such as Randy Barnett, Evan Bernick, and Kurt Lash see a textual error that erased the Clause, and Justice Thomas urges overruling it. Progressives condemn the retreat from Reconstruction but defend the health-law result and note that the dissents seeded \textit{Lochner}.

\section{Key Quotations}
\begin{quote}
``[N]o one can fail to be impressed with the one pervading purpose found in them all, lying at the foundation of each, and without which none of them would have been even suggested; we mean the freedom of the slave race, the security and firm establishment of that freedom, and the protection of the newly-made freeman and citizen from the oppressions of those who had formerly exercised unlimited dominion over him.''
\end{quote}
\noindent --- Justice Miller, for the Court, 83 U.S. at 71.

\end{document}
